\section{Finite-feature Gram identifiability}
\label{sec:gram}
The relevant finite representation is a Gram matrix of \emph{feature
coefficients}, not a sampled data Gram matrix. Let $\Theta\subset\mathbb R^K$
be an open connected parameter box and suppose
\begin{equation}
 f_\theta(x)=C(\theta)h(x)+u(\theta),\quad
 G(\theta)=C(\theta)^T M C(\theta),\label{eq:coefficientgram}
\end{equation}
where $h:D\to\mathbb R^A$ is fixed, $C,u$ are real analytic, and
$M\succ0$ and $\gamma>0$ are fixed. Define $d_\theta^2$ as the
squared $M$-distance of features and $k_\theta=\exp(-\gamma d_\theta^2)$.
We require $\{1,h_1,\ldots,h_A\}$ to be linearly independent as
functions on the \emph{actual} input domain $D$. Equivalently, the
affine hull of $h(D)$ is $\mathbb R^A$.

\begin{theorem}[Coefficient-Gram identifiability]\label{thm:gram}
Under these assumptions, for any $\theta,\eta\in\Theta$,
\begin{equation}
 \begin{gathered}
 k_\theta(x,x')=k_\eta(x,x')\quad\forall x,x'\in D\\
 \Longleftrightarrow\quad G(\theta)=G(\eta).
 \end{gathered}
\end{equation}
For every parameter direction $v$, the kernel directional derivative
vanishes for all input pairs if and only if $D G(\theta)[v]=0$.
Consequently, coordinate $\theta_j$ is globally kernel-inactive if
and only if $\partial_jG\equiv0$, and the generic kernel differential
rank is
\begin{equation}
 r=\operatorname{grank}D_\theta\operatorname{vech}G(\theta)
 \le \min\{K,A(A+1)/2\}.\label{eq:gramrank}
\end{equation}
Here $\operatorname{vech}$ lists the upper-triangular entries once.
\end{theorem}
\begin{proof}
The offset $u$ cancels, so $d_\theta^2=\Delta h^T G(\theta)\Delta h$.
Fix $x_0$ and put $y(x)=h(x)-h(x_0)$. These vectors span $\mathbb R^A$.
If a symmetric $H$ obeys $\Delta h^TH\Delta h=0$ for all pairs, then
$y(x)^THy(x)=0$ and polarization gives $y(x)^THy(x')=0$ for all
$x,x'$. Spanning implies $H=0$. This proves injectivity of the linear
map from $G$ to squared distances. The exponential is injective and
has nonzero derivative, establishing both equivalence and the
directional statement. Integration along coordinate intervals proves
the global inactivity assertion. The finite-dimensional injective
distance map preserves differential rank, giving Eq.~\eqref{eq:gramrank}.
Analyticity makes the maximum rank generic on the connected box.
\end{proof}

Linear independence of $h$ alone is insufficient: $h(x)=(1,x)^T$
has independent components, but $G=\diag(1,1)$ and $G'=\diag(2,1)$
give the same distances. Constant modes must therefore be separated
into $u$. Likewise, formal independent variables cannot replace
dependent coordinates of a restricted data manifold.

For fixed-frequency Pauli circuits, $h$ consists of real sine and
cosine modes with distinct nonzero frequency vectors (one representative
per sign pair). These functions and $1$ are independent if the
transformed data domain contains an open set. The finite-basis representation links all-input identifiability to
coefficient algebra. Its relation to existing spectral methods is
discussed in Appendix~\ref{app:spectral_context}.

\begin{corollary}[Instantaneous kernel-loss gradient space]
\label{cor:gradientspace}
Under Theorem~\ref{thm:gram}, set $g(\theta)=\operatorname{vech}G(\theta)$
and $J_g=D_\theta g$. For a differentiable loss depending on $\theta$
only through finitely many fixed-input kernel evaluations,
\begin{equation}
 \nabla_\theta L\in\operatorname{im}J_g(\theta)^T
       =(\ker J_g(\theta))^\perp.
\end{equation}
This space has dimension $\operatorname{rank}J_g(\theta)\le r$,
with equality on the regular parameter set.
\end{corollary}
\begin{proof}
Each evaluated kernel is a smooth function of $g$, so the loss factors
as $L=\ell\circ g$. The chain rule gives
$\nabla_\theta L=J_g^T\nabla_g\ell$; the orthogonal-complement
identity follows from finite-dimensional linear algebra.
\end{proof}
This is an instantaneous constraint, not confinement to one fixed
$r$-dimensional linear subspace along training. The subspace can rotate
with $\theta$; finite steps, momentum and preconditioning need not stay
in it. Direct parameter regularization can add components outside it.
Nor does this imply $r$ circuit evaluations suffice for a gradient.


\begin{corollary}[Finite decision procedure on an exact subclass]
\label{cor:gramdecision}
Suppose, in addition, the non-data gate angles are rational-affine
functions of $\theta$, all fixed trigonometric constants and metric
entries are in an effective algebraic number field, and an independent
finite data basis has been constructed exactly. Global coordinate
inactivity and the generic differential rank are decidable by finitely
many exact coefficient tests.
\end{corollary}
\begin{proof}
After clearing angle denominators, entries of $G$, their derivatives,
and all Jacobian minors are finite Laurent polynomials in
$e^{i\theta_j/q}$ over an effective algebraic field (adjoining $i$).
Distinct parameter frequencies are independent on an open box.
Thus a derivative or minor vanishes identically exactly when every
coefficient vanishes. Test derivatives for inactivity and minors in
decreasing order for rank. There are finitely many of each.
\end{proof}
This is a termination result, not a polynomial-time guarantee or a
claim that the present pilot implements exhaustive minor enumeration.
Its expansion and minor counts can be prohibitive. Arbitrary analytic
angles, unspecified real constants, and products $\theta_j\phi(x)$
are outside this decision claim. In particular, the ChebyshevRx and
ParamZ examples still require the separate methods below.

\begin{corollary}[Readout-limited dimension bound]\label{cor:readoutbound}
Let $\rho=\min(m,A)$, where $m$ is the number of measured features.
Under Theorem~\ref{thm:gram},
\begin{equation}
 r\le\min\{K,A\rho-\rho(\rho-1)/2\}.
\end{equation}
\end{corollary}
\begin{proof}
The Gram matrix is positive semidefinite of rank at most $\rho$.
The rank-$q$ positive-semidefinite stratum has dimension
$Aq-q(q-1)/2$: factor it by a full-column-rank $A\times q$ matrix,
then quotient by $O(q)$. On the generic constant-rank stratum of
$C(\theta)$, the differential image lies in this stratum's tangent
space. Maximizing over $q\le\rho$ gives the bound.
\end{proof}
This standard rank-stratum bound sharpens Eq.~\eqref{eq:gramrank}
when $m<A$. Its null directions may depend on $\theta$; they need
not be deletable coordinates or executable savings.

\subsection{Isometries and readout dependence}
If $C(\eta)=Q C(\theta)$ with $Q^TMQ=M$, then the kernels agree.
Conversely, equality of their coefficient Grams admits such a $Q$:
whiten by $M^{1/2}$, map corresponding column spans isometrically,
and extend orthogonally to the ambient feature space. Arbitrary
data-independent offsets remain unidentifiable. This describes
reachable-feature equivalence; the all-state Pauli-output
criterion has a stronger universal quantifier.
For a fixed proposed feature transformation $Q$, the exact
reachable-family condition is
$C(\theta)^T(Q^TMQ-M)C(\theta)=0$ for every $\theta$.
It does not generally imply isometry on the linear span of all
these column spaces: cross-parameter bilinear terms are untested.
Appendix~\ref{app:followup} supplies a counterexample.
The Hubregtsen example below is directly represented by a diagonal
$G$ with entries $\sin^2\theta_0+\sin^2\theta_2$ and
$\sin^2\theta_1+\sin^2\theta_3$. A rank deficit need not make any
individual coordinate inactive, nor produce a global minimal chart.
The same example separates $D_\theta\operatorname{vech}G$ from the
coefficient Jacobian $D_\theta\operatorname{vec}C$, whose Gram matrix
of coefficient derivatives is the QNTK-type object built from circuit
harmonic matrices \cite{CHM}. With $h=(\cos x_0,\cos x_1)$, the four
nonzero entries $C_{q,q\bmod2}=-\sin\theta_q$ give
$\operatorname{rank}D_\theta\operatorname{vec}C=4$, one per retained
angle, whereas $\operatorname{rank}D_\theta\operatorname{vech}G=2$,
the number of input groups, at generic $\theta$.

\begin{proposition}[Full Pauli readout and nonmonotonicity]
\label{prop:IC}
For all nonidentity $n$-qubit Pauli words, with their usual
unnormalized convention and unit weights,
\begin{equation}
 d^2=2^n\Tr[(\rho-\rho')^2].
\end{equation}
For pure states this is $2^{n+1}(1-F)$, where $F=\Tr(\rho\rho')$.
For general subsets of Pauli readouts, kernel differential rank is
not monotone under adding measured observables.
\end{proposition}
\begin{proof}
Pauli orthogonality gives the first identity; the trace-zero identity
component contributes nothing. Purity gives the second. For
nonmonotonicity, the certified two-qubit examples in
Table~\ref{tab:dimension} give both directions: at two layers,
ChebyshevRx has $r_Z=6>r_{XYZ}=4$, whereas one-layer ParamZ has
$r_Z=0<r_{XYZ}=2$.
\end{proof}
Adding rows to $C$ adds positive-semidefinite terms to $G$, but this
ordering does not order Jacobian ranks. Full readout is therefore not
generally rank-minimizing. The two-qubit weight-at-most-two readout
is information complete; for mixed states its distance also depends
on the two purities and is not just a function of an overlap.

\subsection{Exact ranks for the Kyriienko settings}
\label{sec:gramkyri}
The frozen one-layer Kyriienko constructors have
\begin{equation}
 U_\theta(x)=E\Bigl(\bigotimes_{q=0}^{n-1}
 R_z(c_q)R_x(b_q)R_z(a_q)\Bigr)D(x),\label{eq:kyriform}
\end{equation}
with $K=3n$. Here $D$ is the fixed-frequency encoding,
$E=\mathrm{CX}_{0,1}$ for $n=2$, and $E$ is the chronological
chain $(\mathrm{CX}_{0,1},\mathrm{CX}_{2,3},\mathrm{CX}_{1,2})$
for $n=4$. Disjoint gates commute to give this form exactly.
Both $Z_0$ and $X_{n-1}$ remain unchanged under conjugation by $E$.
Local XYZ and complete weight-at-most-two readouts with unit weights
are invariant under arbitrary local output rotations. Thus left
rotations about these two axes give independent one-dimensional
kernel gauge orbits. On a regular Euler chart ($\sin b_q\ne0$),
the Euler map is locally invertible, so the two source-parameter
tangent directions are independent and $r\le3n-2$.
Analyticity extends the vanishing higher minors from this open set.

For a matching lower bound we compute real data-Fourier coefficients
and their parameter derivatives using exact Gaussian rational
arithmetic at
\begin{equation}
 \theta_j=2\arctan t_j,\qquad t_j=(j+2)/(j+5).
\end{equation}
All sines and cosines are rational. Nonzero exact Gram-Jacobian minors
of orders $4,10,10$ prove $r\ge4,10,10$ for the two-qubit XYZ,
four-qubit XYZ, and four-qubit weight-at-most-two settings. Together
with the gauge upper bounds, these minors fix all three ranks exactly.
These are three settings of the existing census, not three new
independent families. Full rational matrices and determinants are
retained in the accompanying audit; no floating rank threshold is used.

The resulting ranks are exactly $4,10,10$, closing all 54 settings
across the two dimension studies. The second gauge is a
parameter-dependent Euler direction, not another globally dead
source coordinate. These closures consequently do not change the
reported inactive-coordinate counts or establish gate or shot savings.
Implementation and independent source checks are given in
Appendix~\ref{app:gram}.
